\section{Related Work}
\label{sec:related}

\paragraph{E-graphs and congruence closure.}
The e-graph data structure was introduced by Greg Nelson in his 1980 PhD
thesis~\cite{nelson1980}, alongside the congruence-closure decision procedure
developed jointly with Derek Oppen~\cite{nelsonoppen}. The e-graph was originally
designed as the core data structure of the simplifier in the Stanford Pascal
Verifier, where it represented equivalence classes of ground terms under the
axioms of equality and uninterpreted functions. Modern SMT solvers and
equality saturation engines inherit this design.

\paragraph{E-matching in theorem provers.}
The idea of \emph{matching modulo the equivalence stored in the e-graph} --- now
called e-matching --- already appears in Nelson's thesis~\cite{nelson1980} . The
technique was refined and used in the ESC line of program checkers, but its
first detailed algorithmic exposition is the canonical Simplify paper of
Detlefs, Nelson, and Saxe~\cite{simplify}. Our work targets the same problem ---
finding all matches of a pattern modulo an e-graph --- but compares Simplify's
recursive top-down matcher with a relational e-matching algorithm.

\paragraph{Indexing for SMT e-matching.}
de Moura and Bj{\o}rner introduced the \emph{e-matching code tree}, a
compiled abstract matching machine that shares prefix work across many
patterns simultaneously, and the \emph{inverted path index}, which makes
e-matching incremental as the e-graph grows~\cite{ematchingz3}. A companion
paper~\cite{ematchingfunprofit} surveys the practical lessons from deploying
these techniques in Z3 and motivates index-based matching as the workhorse of
modern SMT quantifier instantiation. Code trees are the closest prior work to
our approach in spirit: both compile patterns to a more efficient
representation than naive top-down traversal. The key difference is that code
trees still drive matching from the pattern root, whereas relational
e-matching is free to start from whichever relation is smallest.

\paragraph{Equality saturation.}
Equality saturation, introduced by Tate et al.~\cite{eqsat}, repurposes the
e-graph as a compiler intermediate representation: rewrite rules are applied to
saturation and an optimal program is then extracted, sidestepping the
phase-ordering problem in traditional optimizers. The egg library~\cite{egg}
made this approach widely practical by introducing deferred rebuilding of
congruence invariants and e-class analyses. Relational
e-matching~\cite{relationalematching} was developed in this lineage and
implemented in egglog~\cite{egglog}, where the absence of backtracking makes
maintaining the relational indices straightforward. Our contribution is to push
these ideas back into the SMT setting, where backtracking is unavoidable and the
index structures must be made backtrackable. Another difference is that egglog 
employs a worst-case optimal join algorithm~\cite{wcoj}, whereas our current implementation uses a
left-deep binary join. We leave the implementation of a worst-case optimal
join for future work.
